\section{Introduction}
\label{sec:intro}

Distributed cyber-physical systems (dCPS) are increasingly deployed on a heterogeneous
compute continuum that spans sensing devices, edge nodes, gateways and cloud
infrastructure. Designing such a system is not a single decision but a joint one: how
many edge nodes and gateways to deploy, which hardware class each of them should be,
where each application stage is executed, how far the ingest stage is replicated, and
which class of communication link connects the tiers. These decisions interact. Adding
edge nodes reduces queueing delay but raises deployment cost and static energy; moving
analytics to the cloud removes load from the edge tier but adds a wide-area transfer to
every request; replication improves ingest throughput but multiplies both compute load
and synchronisation traffic.

Because the decisions are discrete and interacting, the number of candidate
architectures grows as the product of the decision domains and quickly reaches $10^{5}$
to $10^{7}$ even for a modest set of design variables. At the same time, evaluating a
single candidate is expensive: a recent industrial dCPS exploration methodology reports
that ``evaluating each design point takes under a minute'' with a discrete-event
simulator, and its authors explicitly leave the integration of a search strategy out of
scope~\cite{saadatmand2025compdse}. A survey of the field concludes that ``efficient and
scalable DSE technology for dCPS is more or less non-existing and constitutes a largely
unchartered research area'' and identifies the cost of a design-point evaluation as a
defining difficulty of the distributed setting~\cite{herget2022dse}. The combination of
a large space and an expensive evaluator is precisely the regime in which the number of
\emph{full evaluations}, rather than the raw size of the space, determines whether
design-space exploration (DSE) is practical at all.

Two families of approaches dominate. Exhaustive or model-based enumeration evaluates
every candidate and therefore yields the exact Pareto set, but its cost scales with the
size of the space. Population-based metaheuristics such as NSGA-II~\cite{deb2002nsga2}
operate under a fixed evaluation budget and return an approximation with no guarantee on
which trade-offs have been missed; as we show in Section~\ref{sec:results}, a
budget-matched NSGA-II run needs $10^{3}$ evaluations to recover half of the exact front
on a space of $5\times10^{4}$ candidates. A third, smaller family performs \emph{exact}
multi-objective exploration by pruning partial designs, either symbolically with
constraint solvers~\cite{neubauer2018exact,haubelt2018asp} or through bound-based Pareto
pruning in the broader optimisation literature~\cite{paretopruning2022,mandow2010namoa}.
Methods in this family share one structural assumption: the objectives must be
\emph{assignment monotonous}, i.e.\ each additional decision can only add to the
objective, so that a partial design already carries a valid lower bound on all its
completions.

That assumption does not hold for the objective that matters most in a distributed CPS.
End-to-end latency contains a waiting component that depends on the utilisation of a
tier, and utilisation is a function of several decisions at once --- the number of nodes
in the tier, their class, the replication factor and the placement policy. No individual
assignment increases latency monotonically, so the standard ``cost so far plus optimistic
remainder'' bound is unavailable. This is the gap the present paper addresses.

\paragraph{Contributions}
\begin{enumerate}
\item A formal design-space model for heterogeneous edge--gateway--cloud CPS
  architectures (Section~\ref{sec:model}) with analytical models of latency including
  contention, energy, normalised deployment cost, network volume, and six classes of
  feasibility constraints.
\item \textbf{Admissible bounds for contention-dependent objectives}
  (Section~\ref{sec:method}). We decompose each objective into non-negative terms,
  bound every term independently over the still-reachable domains, and exploit a
  decision ordering under which the determinants of tier utilisation are fixed early, so
  that the waiting term becomes \emph{exactly} computable on a partial design rather
  than being relaxed to zero. Lemmas~\ref{lem:relax} and~\ref{lem:exact} establish
  admissibility.
\item \textbf{HCA-DSE}, a hierarchical exploration method that combines structural
  well-formedness, feasibility pruning and bound-based dominance pruning with a final
  exact evaluation of the survivors, together with proofs
  (Propositions~\ref{prop:feas}--\ref{prop:struct}, Theorem~\ref{thm:exact}) that the
  exact Pareto front is preserved.
\item \textbf{An evaluation-cost analysis} that answers when such pruning is worth its
  own overhead, including a measured break-even evaluator cost and a crossover curve
  (Section~\ref{sec:results}).
\item \textbf{An empirical validation} on a containerised edge--gateway--cloud testbed
  with per-node link shaping: 36 architectures, 108 measured runs, reported both as
  prediction error and as the agreement between the architecture ordering produced by
  the model and the ordering observed in the measurements.
\item A reproducibility package in which every figure and every numeric table of this
  paper is generated from the raw result files by a single script, and in which the
  correctness claims are backed by exhaustive machine-checked tests.
\end{enumerate}

The remainder of the paper is organised as follows. Section~\ref{sec:related} positions
the work. Section~\ref{sec:model} defines the design space and the evaluation model.
Section~\ref{sec:method} presents HCA-DSE and its correctness.
Section~\ref{sec:methodology} describes the experimental methodology,
Section~\ref{sec:results} the results. Section~\ref{sec:discussion} discusses when the
method pays off and when it does not, Section~\ref{sec:threats} the threats to validity,
and Section~\ref{sec:conclusion} concludes.
